\documentclass[11pt,a4paper]{article}

\usepackage[utf8]{inputenc}
\usepackage[T1]{fontenc}
\usepackage{amsmath,amssymb,amsfonts}
\usepackage{graphicx}
\usepackage{booktabs}
\usepackage{hyperref}
\usepackage{xcolor}
\usepackage{microtype}
\usepackage[margin=1in]{geometry}
\usepackage{natbib}
\usepackage{caption}
\usepackage{enumitem}

\hypersetup{colorlinks=true,linkcolor=blue,citecolor=blue,urlcolor=blue}

\title{Bankai: Ultra-Sparse Adaptation of 1-Bit LLMs via XOR Patches}

\author{
  Nikshep Saravanan \\
  \texttt{github.com/nikshepsvn/bankai}
}

\date{April 2, 2026}

\begin{document}
\maketitle

% ============================================================
\begin{abstract}
True 1-bit large language models have no post-training adaptation method.
LoRA, fine-tuning, and quantization-aware training all require continuous
weights or gradients that binary models lack. We introduce \textbf{Bankai},
the first post-training adaptation method for true 1-bit LLMs, using
bitwise XOR operations on binary weights. Bankai patches are sparse XOR
bitmasks that modify model weights in-place with a single bitwise operation,
incur zero inference overhead, and compress to around one kilobyte.

We validate on Bonsai~8B (PrismML, 2026), a true 1-bit, 8.2~billion
parameter language model. Through eleven experiments, we demonstrate that:
(1)~binary MLP weights exhibit massive redundancy;
(2)~scale-guided bit flips produce $3.88\times$ more behavioral impact
than random flips;
(3--4)~greedy search yields patches that correct specific calculus failures
in free generation;
(5)~patches trained on few probes memorize rather than generalize; and
(6)~training on diverse probe variations produces patches that
\textbf{generalize to held-out prompts}---fixing 4~of~17 problems
the base model gets wrong (23.5\%) with zero breakage on the 13 it
already solves.

\medskip
\noindent\textbf{Correction to this version.} The Experiment~6 figure quoted
above (4~of~17, 23.5\%) is superseded. An audit prompted by an independent
reproduction found three defects in the probe measurement; under corrected
scoring the same patch fixes \textbf{2}~held-out problems with zero breakage,
confirmed at the generation level rather than by logit gaps. The central claim
is unaffected. See Section~\ref{sec:errata}.
\end{abstract}

% ============================================================
\section{Introduction}

The emergence of true 1-bit large language models---where every weight is
a single bit---represents a fundamental shift in model efficiency.
Bonsai~8B~\citep{prismml2026} stores 8.2~billion parameters in 1.15~GB,
runs at 44~tokens/second on an iPhone, and achieves benchmark scores
competitive with full-precision 8B~models.

However, these models present a novel challenge: \textbf{they cannot be
adapted after deployment}. Every existing method for post-training
adaptation requires continuous-valued operations:

\begin{itemize}[nosep]
  \item \textbf{LoRA}~\citep{hu2021lora} adds low-rank weight deltas in
    float space.
  \item \textbf{Fine-tuning} backpropagates continuous gradients through
    differentiable weights.
  \item \textbf{Quantization-aware training} adjusts weights during
    pre-training, not after deployment.
\end{itemize}

None of these apply to a model whose weights are single bits. Once a 1-bit
model is deployed, it is frozen.

We introduce \textbf{Bankai}, the first post-training adaptation method
for true 1-bit LLMs. The key insight: because weights are bits, the
``diff'' between two behavioral states is a bitwise XOR mask. Where the
mask has a~1, a weight flipped. Where it has a~0, nothing changed. For
nearby behavioral variants, this mask is sparse---enabling patches that
modify $0.007\%$ of weights (${\sim}1$~KB) with zero inference overhead,
microsecond application, and exact reversibility.

We validate Bankai on Bonsai~8B through six experiments, progressing from
mechanism validation (Experiments~1--2), through targeted correction
(Experiments~3--4), to generalization testing (Experiments~5--6). We find
that patches trained on 6~probes memorize specific prompts, but patches
trained on 60~diverse probes generalize to held-out problems---fixing
4~of~17 incorrect answers with zero collateral damage.

\subsection{Contributions}

\begin{enumerate}[nosep]
  \item \textbf{The first post-training adaptation method for true 1-bit
    LLMs.} Existing methods (LoRA, fine-tuning, QAT) require continuous
    weights or gradients and are inapplicable to binary architectures.
  \item \textbf{Scale-guided targeting.} Per-group FP16 scale factors
    predict which weight regions are most behaviorally sensitive, enabling
    $3.88\times$ more efficient search than uniform random sampling.
  \item \textbf{Layer-level functional specialization in 1-bit
    architectures.} Early layers (1--4) and the penultimate layer~(34)
    are load-bearing, while middle layers (17--21) contribute minimally.
  \item \textbf{A characterization of generalization.} Patches trained on
    few probes memorize; diverse training produces generalization. This
    is the first empirical characterization of how XOR patches learn.
  \item \textbf{An open toolkit and patch format} for reproducible
    creation, application, and evaluation of XOR patches.
\end{enumerate}

% ============================================================
\section{Background}

\subsection{True 1-Bit vs.\ Ternary Models}

The distinction between true 1-bit and ternary models is critical.
BitNet~b1.58~\citep{wang2023bitnet,ma2025bitnet} uses ternary weights
$\{-1, 0, +1\}$, requiring 2~bits per weight in storage. Despite the
``1-bit'' branding, XOR on 2-bit ternary encodings produces invalid
states: $\text{XOR}(01, 10) = 11$, which has no valid ternary mapping.

Bonsai~8B~\citep{prismml2026} is a true 1-bit model where each weight
is a single bit: $0 \to -\text{scale}$, $1 \to +\text{scale}$. Weights
are packed as \texttt{uint32} arrays with per-group FP16 scale factors
(group size~128). This makes bitwise XOR a semantically valid operation.
As of April~2026, Bonsai~8B is the only production-quality true 1-bit LLM.

\subsection{Deployment Implications}

The properties of XOR patches---kilobyte-scale size, microsecond
application, zero inference overhead, instant reversibility---enable
a deployment model impossible with existing adaptation methods.
A library of domain patches (math, code, medical, legal), each
${\sim}1$~KB, can be hot-swapped at inference time on a phone.
The base model stays at 1.15~GB; a thousand patches add 1~MB.
LoRA adapters are ${\sim}100$~MB each, add per-token latency,
and require weight reloading to swap.

% ============================================================
\section{Related Work}

\textbf{Low-rank adaptation.}
LoRA~\citep{hu2021lora} and its variants (QLoRA~\citep{dettmers2023qlora};
DoRA) reduce fine-tuning cost via low-rank weight deltas. These require
continuous weights and are not applicable to true 1-bit architectures.

\textbf{Compact adaptation.}
RECAST~\citep{xu2024recast} reduces task-specific parameters to fewer
than~50 via weight reconstruction on continuous-weight models.

\textbf{Binary neural networks.}
XOR-Net~\citep{bulat2020xornet} and XNOR-Net~\citep{rastegari2016xnor}
use binary operations for efficient forward passes, not post-training
behavioral modification.

\textbf{1-bit and sub-1-bit LLMs.}
BitNet~\citep{wang2023bitnet} introduced 1.58-bit LLM training.
STBLLM~\citep{dong2024stbllm} pushed compression below 1-bit and
observed that some weights in binarized LLMs can be randomly flipped
without significant performance degradation---a finding our
Experiment~1 independently confirms.

\textbf{Bit-flip attacks.}
\citet{rakin2019bitflip} showed that ${\sim}20$ targeted bit flips can
catastrophically degrade a quantized DNN. Subsequent work
(T-BFA~\citep{bai2021tbfa}; Versatile Weight Attack) refined targeted
attacks. This literature establishes that small numbers of bit flips
produce outsized behavioral effects---but focuses exclusively on
adversarial degradation. \textbf{Bankai inverts this mechanism:}
constructive bit flips for targeted capability improvement.

\textbf{Model editing.}
ROME~\citep{meng2022rome} and MEMIT~\citep{meng2023memit} perform
targeted factual edits on continuous-weight models. Bankai shares
the goal of minimal intervention but operates on binary weights for
behavioral (not factual) modification.

% ============================================================
\section{Methodology}

All experiments use Bonsai~8B (\texttt{prism-ml/Bonsai-8B-mlx-1bit}),
a true 1-bit, 8.2B~parameter model based on the Qwen3 architecture
(36~layers, 4096 hidden dim, GQA with 32/8~heads). Weights are packed
as \texttt{uint32} arrays with 1-bit group quantization
(\texttt{group\_size=128}). Experiments were run on Apple~M3 (24~GB).

\subsection{Probe Metrics}

Experiments~1--11 use the \texttt{token\_gap} metric described first; it has
defects documented in Section~\ref{sec:errata}. Experiment~12 uses
\texttt{seq\_logprob} for search and greedy generation for reporting.

\paragraph{\texttt{token\_gap} (Experiments 1--11).}
Pairs of (correct\_token, wrong\_token) following a deterministic prompt. The
logit gap $G = \text{logit}(\text{correct}) - \text{logit}(\text{wrong})$
is a single-forward-pass measurement: positive means the model prefers
the correct answer. Answer strings are reduced to a single token id by keeping
their last subtoken. We used this as a fast proxy for behavioral change,
noting that logit improvement does not always translate to generation
improvement (see Section~\ref{sec:limitations}).

\paragraph{\texttt{seq\_logprob} (Experiment 12).}
The gap is the difference in summed teacher-forced log-probability between the
full correct answer string and a plausible per-probe distractor, with
continuations tokenized in the context of their prompt. Multi-token answers,
leading spaces, and same-suffix collisions are all handled correctly. This is
used as search fitness only.

\paragraph{Greedy generation (Experiment 12, headline).}
What the model actually emits under greedy decoding, compared against the
expected answer and any accepted alternate renderings. It depends on no
distractor token, so it cannot be satisfied by shifting probability between two
hardcoded ids. Logit gaps are a search signal; generation is the result.

\subsection{Backends: MLX and GGUF/CUDA}

Bankai has two independent backend implementations:

\begin{itemize}[nosep]
  \item \textbf{MLX} (Apple Silicon): uses PrismML's MLX fork with
    1-bit kernels. Implements \texttt{flip\_row} as a bitwise XOR on
    packed \texttt{uint32} weight tensors, followed by
    \texttt{mx.eval()} to propagate the modification.
  \item \textbf{GGUF/CUDA} (NVIDIA GPUs): uses a custom C++ tool,
    \texttt{bankai\_eval.cpp}, built against PrismML's llama.cpp fork.
    The tool loads Bonsai~8B once and exposes a line protocol on
    stdin/stdout (\texttt{PROBE}, \texttt{TOKENIZE}, \texttt{FLIP\_ROW},
    \texttt{FLIP\_GROUP}, \texttt{SCALES}, \texttt{NUM\_ROWS}).
    Weight modifications happen entirely in GPU memory via
    \texttt{ggml\_backend\_tensor\_get}/\texttt{set} on the
    Q1\_0\_g128 tensor data --- no file reloading, no subprocess
    restart. The Python \texttt{GGUFBackend} wraps this subprocess
    and pipelines probe commands through the stdin buffer, so a full
    multi-probe iteration costs one write, one flush, and one batched
    read. This achieves approximately $24\times$ speedup over the MLX
    path on Apple M3: a 300-iteration search runs in $\sim$3~min
    on Modal L40S vs $\sim$67~min on M3 MacBook Air.
\end{itemize}

Patch files (JSON lists of \texttt{(layer, proj, row[, group])}
tuples) are portable between the two backends as long as the tensor
names match. The MLX and GGUF versions of Bonsai differ in
quantization format, however (MLX uses scale + bias per 128-weight
group at 1.25~bpw; GGUF uses scale-only at 1.125~bpw), so identical
patches can produce different logit gaps on the two formats ---
this is discussed in Experiment~9.

\subsection{Experiment~1: Random Bit Flip Robustness}

We flip $N$ random bits across all MLP weight tensors (5.4B~bits total)
and measure perplexity on a fixed 5-sentence eval set, for
$N \in \{100, 1\text{K}, 10\text{K}, 50\text{K}, 100\text{K}, 500\text{K}\}$
under four strategies: random across all layers, random in layers~16--24,
scale-guided medium (25th--75th percentile), and scale-guided high
(top~25\%).

\subsection{Experiment~2: Structured Flips and Layer Specialization}

For each of 36~layers, we flip all bits in the entire MLP and measure
logit gaps on 8~probes spanning arithmetic, code, geography, and
science. We then test row-level flips at varying counts (1--1024~rows)
in the most impactful layer and compare high-scale rows vs.\ random
rows at 64~rows.

\subsection{Experiment~3: Greedy Patch Search (Arithmetic)}

Greedy hill climbing over row-level flips in layers $[1, 2, 3, 4, 34]$
(selected from Experiment~2). Each iteration: sample a row weighted by
scale magnitude, flip all 4,096~bits, measure fitness, keep if improved,
revert if not. 200~iterations with 6~target probes (arithmetic) and
4~control probes (knowledge).

\textbf{Fitness function:}
\begin{equation}
  f = \frac{1}{|T|}\sum_{t \in T} (g_t - g_t^0)
    - \lambda \cdot \frac{1}{|C|}\sum_{c \in C} \max(0, g_c^0 - g_c)
\end{equation}
where $g_t, g_c$ are target and control logit gaps, superscript~$0$
denotes baselines, and $\lambda = 2.0$ penalizes control degradation.

\subsection{Experiment~4: Calculus Patch}

Same search targeting 6~calculus probes (polynomial differentiation,
second derivatives, integrals, primality, trigonometry, exponential
derivatives) where the base model fails deterministically (verified
across 5~runs). 300~iterations with screening optimization.

\subsection{Experiment~5: Variation Testing}

We generate 15~novel variations per probe category (90~total) and
measure sign flips (wrong$\to$right vs.\ right$\to$wrong) on the
Experiment~4 patch applied to these never-seen probes.

\subsection{Experiment~6: Generalization-Optimized Search}

Train on 60~probes (10~per category with diverse phrasings), hold out
30~probes (5~per category) for validation. Mean fitness, 300~iterations.
Search time scales roughly linearly with probe count:
6~probes $\to$ 13~min, 60~probes $\to$ 67~min.

\subsection{Experiment~7: Patch Stacking}

Apply the Experiment~3 math patch (72~flips) and Experiment~4 calculus
patch (70~flips) simultaneously via sequential XOR. Verify order
independence and reversibility. Test both math and calculus probes.

\subsection{Experiment~8: GSM8K Safety Check}

Run 50~GSM8K word problems with full generation (400~tokens) and answer
extraction. Compare accuracy with and without the Experiment~6
generalized patch to check for collateral damage on general math
reasoning.

\subsection{Experiment~9: GGUF/CUDA Backend (Beefier Search)}

Port Bankai to PrismML's llama.cpp fork (Q1\_0\_g128 CUDA kernels) via
a custom C++ tool, \texttt{bankai\_eval}, that performs in-GPU weight
manipulation through \texttt{ggml\_backend\_tensor\_get}/\texttt{set}.
Run the Experiment~6 search on Modal L40S with an expanded layer set
$[1, 2, 3, 4, 5, 6, 10, 34]$ (adding layers 5, 6, 10 which the
Experiment~2 layer-impact map identified as calculus-sensitive) and
800 iterations instead of 300.

\subsection{Experiment~10: Attention Projection Search}

Same configuration as Experiment~9 but with
\texttt{search\_projs = [gate, up, q, k, v, o]} --- the search now
considers attention Q/K/V/O projections in addition to MLP. First
exploration of attention-projection XOR flips on a 1-bit LLM.

\subsection{Experiment~11: Per-Group Granularity Search}

Replace row-level flipping with group-level: each ``flip'' XORs only
the 128 sign bits of one 18-byte Q1\_0\_g128 block instead of all
4,096 sign bits of a full row. Search space grows $32\times$ (from
$\sim$131K rows to $\sim$6.3M groups). 1500 iterations, same 8 layers
and 2 MLP projections.

% ============================================================
\section{Results}

\subsection{Experiment~1: Random Flips Are Absorbed}

\begin{table}[h]
\centering
\caption{Perplexity change under random bit flips in MLP weights.}
\begin{tabular}{@{}lll@{}}
\toprule
Flips & \% of MLP bits & Max Perplexity $\Delta$ \\
\midrule
100     & 0.000002\% & $< 0.01$ \\
10,000  & 0.0002\%   & $< 0.01$ \\
100,000 & 0.002\%    & $< 0.02$ \\
500,000 & 0.009\%    & $< 0.08$ \\
\bottomrule
\end{tabular}
\end{table}

The model is remarkably robust to random perturbation in MLP weights,
implying massive redundancy. (Attention weights, embeddings, and the
LM~head were not tested.)

\subsection{Experiment~2: Layer Impact and Scale-Guided Targeting}

\begin{table}[h]
\centering
\caption{Average absolute logit gap change ($|\Delta\text{gap}|$) when
  flipping the entire MLP of each layer range. 8~probes across 4~domains.}
\begin{tabular}{@{}lll@{}}
\toprule
Layer range & Avg $|\Delta\text{gap}|$ & Interpretation \\
\midrule
0--4 (early)        & 3.2--7.2 & High impact \\
5--16 (middle)      & 0.7--3.0 & Moderate, decreasing \\
17--21 (deep mid)   & 0.7--1.6 & \textbf{Lowest impact} \\
22--33 (late)       & 1.6--3.4 & Moderate, increasing \\
34 (penultimate)    & \textbf{9.0} & \textbf{Highest impact} \\
35 (final)          & 3.2 & High \\
\bottomrule
\end{tabular}
\end{table}

\textbf{Scale-guided targeting:} At 64~rows in layer~34, high-scale
rows produce $3.88\times$ more behavioral impact than random rows
(avg $|\Delta\text{gap}|$: 0.459 vs.\ 0.118). Layer effects appear
domain-specific across our probe set---layer~34 shifts physics by
$+13.8$ while shifting code by $-15.9$.

\subsection{Experiment~3: Arithmetic Patch}

200~iterations, 72~accepted flips, 7.5~minutes on M3. The \texttt{2+2}
probe flips from incorrect (prefers~5) to correct (prefers~4) in both
logit gap and free generation. The \texttt{7$\times$8} probe improves
at the logit level (prefers~56 over~54) but free generation produces~64
($=8\times8$)---a different wrong answer, illustrating the gap between
probe-based and generation-based evaluation.

\textbf{Patch:} 72~flips, 294,912~bits modified (0.005\%), 864~bytes.
No degradation on 4~measured control probes.

\subsection{Experiment~4: Calculus Patch}

300~iterations, 70~accepted flips, ${\sim}13$~minutes on M3. Two
deterministic failures corrected in free generation:

\begin{itemize}[nosep]
  \item \texttt{d/dx [x\^{}4 + 3x\^{}2] = 0}
    $\to$ \texttt{4x\^{}3 + 6x} \checkmark
  \item \texttt{The second derivative of x\^{}4 is 12x\^{}3}
    $\to$ \texttt{12x\^{}2} \checkmark
\end{itemize}

\textbf{Prompt sensitivity finding:} The primality probe
(\texttt{Is 97 prime?}) is sensitive to a trailing space---the base
model's answer flips from ``No'' to ``Yes'' without any patch. This
illustrates the fragility of single-token probes and motivates
diverse training sets (Experiment~6).

\textbf{Patch:} 70~flips, 286,720~bits (0.005\%), 840~bytes.
No degradation on 5~measured control probes.

\subsection{Experiment~5: Variation Testing (Negative Result)}

\noindent\fbox{\begin{minipage}{0.96\linewidth}
\textbf{Superseded---and understated.} 7 of these 90 variation probes are dead
under \texttt{token\_gap} (Section~\ref{sec:errata}). Rerun under corrected
scoring (Section~\ref{sec:exp13}), this patch breaks \textbf{15} probes rather
than 5, taking accuracy from 49/90 to 36/90. The conclusion below is correct and
in fact too generous. Numbers left unmodified as the published record.
\end{minipage}}


The Experiment~4 patch was tested on 90~novel variations (15~per
category). Sign-flip analysis:

\begin{table}[h]
\centering
\caption{Experiment~4 patch applied to 90~never-seen probes.}
\begin{tabular}{@{}lr@{}}
\toprule
Metric & Count \\
\midrule
Fixed (wrong $\to$ right) & 2 \\
Broke (right $\to$ wrong) & 5 \\
Stayed right & 57 \\
Stayed wrong & 26 \\
\midrule
Net sign flips & $-3$ \\
Base accuracy & 62/90 (68.9\%) \\
Patched accuracy & 59/90 (65.6\%) \\
\bottomrule
\end{tabular}
\end{table}

\textbf{Finding:} The 6-probe patch memorizes specific prompts rather
than shifting capabilities. All 5~broken probes had borderline baseline
gaps ($< 1.2$).

\subsection{Experiment~6: Generalization-Optimized Search}

\noindent\fbox{\begin{minipage}{0.96\linewidth}
\textbf{Superseded.} The fix counts in this subsection are inflated by probe
measurement defects (Section~\ref{sec:errata}). Under corrected scoring the same
patch fixes \textbf{2} held-out problems, not 4, with zero breakage; three of the
four ``fixes'' were already correct at baseline. Numbers are left unmodified as
the published record.
\end{minipage}}


Motivated by Experiment~5's memorization finding, we trained on
60~probes (10~per category with diverse phrasings) and validated on
30~held-out probes.

\begin{table}[h]
\centering
\caption{Sign-flip analysis on training and held-out validation sets.}
\begin{tabular}{@{}lrr@{}}
\toprule
Metric & Training (60) & Validation (30) \\
\midrule
Fixed (wrong $\to$ right) & 4 & \textbf{4} \\
Broke (right $\to$ wrong) & 0 & \textbf{0} \\
Stayed right & 44 & 13 \\
Stayed wrong & 12 & 13 \\
\midrule
Accuracy & 44/60 $\to$ 48/60 & 13/30 $\to$ \textbf{17/30} \\
\bottomrule
\end{tabular}
\end{table}

Validation fixes span 4~categories: polynomial derivatives, second
derivatives, integrals, and primality. Trig and exponential derivative
categories saw zero fixes---a calculus-specific layer impact analysis
identified layers~5, 6, and~10 as high-impact for calculus but not
included in the search set, which may explain the gap.

The base model gets 17~of~30 validation probes wrong; the patch fixes
4~of~those~17 (23.5\%) while breaking none of the 13 it already solves.

\textbf{Patch:} 93~flips, 380,928~bits (0.007\%), 1,116~bytes.
No degradation on 5~measured control probes.

\subsection{Experiment~7: Patch Stacking}

The Experiment~3 math patch (72~flips) and Experiment~4 calculus patch
(70~flips) have zero row overlap, yielding 142~combined flips.
Stacking is mechanically correct: order-independent (math+calc gives
identical results to calc+math) and perfectly reversible (zero drift
after removal).

However, behavioral composition shows interference. The math patch
damages \texttt{poly\_deriv} ($-1.59$) while the calculus patch improves
it ($+2.77$); stacked, they partially cancel ($+0.19$). Overall:
individual patches each fix~2 and break~1 probe; the stacked patch
fixes~1 and breaks~0.

\textbf{Finding:} Patches compose algebraically but interfere
behaviorally. Stacking is safer (fewer breakages) but less effective
(fewer fixes) than individual patches. Joint optimization---searching
for flips that improve both domains simultaneously---would likely
outperform naive stacking.

\subsection{Experiment~8: GSM8K Safety Check}

50~GSM8K word problems, generation with answer extraction:

\begin{table}[h]
\centering
\caption{GSM8K accuracy with and without patch (50~problems).}
\begin{tabular}{@{}lrr@{}}
\toprule
& Without patch & With patch \\
\midrule
Correct & 11/50 & 14/50 \\
Accuracy & 22.0\% & 28.0\% \\
\bottomrule
\end{tabular}
\end{table}

No degradation detected. The patch slightly improved accuracy (+3~problems),
likely within noise for $n=50$ but directionally positive. Our GSM8K accuracy (22\%) is below Bonsai's reported benchmark (88\%)
due to evaluation harness differences; the relative comparison is the
meaningful signal.

\subsection{Experiment~9: GGUF/CUDA Backend}

\textbf{Search:} 800 iterations on Modal L40S, 8 layers
$[1,2,3,4,5,6,10,34]$ $\times$ 2 MLP projections. 424~seconds total ---
approximately $9.5\times$ faster than the original M3 run (4018~s).

\textbf{Patch:} 73 flips (54 screened out), 876~bytes. Layers 5, 6,
and 10 all contributed accepted flips, confirming the
Experiment~2 layer-impact map's prediction.

\begin{table}[h]
\centering
\caption{Experiment~9 results (Modal L40S, GGUF backend).}
\begin{tabular}{@{}lrr@{}}
\toprule
Set & Baseline & Patched \\
\midrule
Training (60 probes) & 43/60 & 42/60 \\
\textbf{Validation (30 held-out)} & \textbf{13/30} & \textbf{16/30} \\
Validation fixed & --- & \textbf{3} \\
Validation broke & --- & \textbf{0} \\
\bottomrule
\end{tabular}
\end{table}

\textbf{Finding:} The GGUF format (scale-only Q1\_0\_g128, 1.125~bpw)
is strictly less expressive than MLX (scale+bias, 1.25~bpw), and
baseline logit gaps differ between formats (for example, Paris-vs-London
is $+6.50$ on MLX vs $+3.54$ on GGUF for identical tokens). Despite the
format difference, the GGUF search with expanded layers recovers
$3/4$ of the MLX validation result (3~fixed vs 4) while running in a
fraction of the time. This establishes the GGUF/CUDA path as a viable
production route.

\subsection{Experiment~10: Attention Projection Search}

\textbf{Search:} Same as Experiment~9 with six projections per layer
(gate, up, q, k, v, o). 800~iterations, 433~seconds.

\textbf{Patch:} 85 flips total --- 69 MLP (gate: 37, up: 32) and
\textbf{16 attention} (q: 7, o: 4, k: 3, v: 2). All four attention
projection types received accepted flips, confirming the mechanism
is mechanically valid on attention weights.

\begin{table}[h]
\centering
\caption{Experiment~10 vs Experiment~9 --- adding attention to the search.}
\begin{tabular}{@{}lrr@{}}
\toprule
 & Exp~9 (MLP only) & Exp~10 (+ attention) \\
\midrule
Training fixed/broke & 2 / 3 & 2 / 1 \\
\textbf{Validation fixed} & \textbf{3} & \textbf{0} \\
Validation broke & 0 & 0 \\
\bottomrule
\end{tabular}
\end{table}

\textbf{Finding (negative result):} Attention flips are mechanically
functional but \textbf{hurt generalization}. Adding the attention
search space to Experiment~9 eliminated all three held-out validation
fixes. Every probe category that had improved in Experiment~9
regressed to zero fixes in Experiment~10.

\textbf{Interpretation:} Attention weights encode context-specific
routing patterns. Flipping rows of attention projections finds
modifications that improve training-set fitness but depend on exact
prompt structure and don't transfer to novel phrasings. MLP weights,
by contrast, encode more abstract representations whose modifications
generalize. \textbf{For behavioral patching of 1-bit LLMs, MLP is the
right target; attention is too context-bound.} This is, to our
knowledge, the first empirical characterization of attention XOR
flips on a 1-bit model.

\subsection{Experiment~11: Per-Group Granularity Search}

\textbf{Search:} 1500 iterations, \texttt{granularity="group"},
755~seconds. Candidate space: 6.3~million group-flips (vs $\sim$131K
row-flips in Experiment~9).

\textbf{Patch:} \textbf{Only 6 flips accepted} (of 1500 iterations,
$\sim$0.4\% accept rate). 331 screened out. Max fitness reached:
$+0.0067$ (vs $+0.32$ for the row-level beefy run in Experiment~9 ---
two orders of magnitude smaller).

\begin{table}[h]
\centering
\caption{Experiment~11: per-group flips produce no measurable behavior change.}
\begin{tabular}{@{}lrr@{}}
\toprule
Set & Baseline & Patched \\
\midrule
Training (60 probes) & 43/60 & 43/60 \\
Validation (30 held-out) & 13/30 & 13/30 \\
Fixed / broke & --- & 0 / 0 \\
\bottomrule
\end{tabular}
\end{table}

\textbf{Finding (negative result):} Per-group flipping is too
fine-grained for the current mean-fitness search. Each 128-bit flip
produces probe gap changes on the order of $0.001$--$0.01$ ---
approximately $10\times$ smaller than row-level flips --- and the
control-degradation penalty ($\lambda=2.0$) dominates at this scale.
The search correctly rejected $\sim$99.6\% of candidates, and the
six accepted flips produced no measurable behavioral change.

\textbf{Interpretation:} Finer granularity is not free. It requires
either (a)~a more sensitive fitness function (log-probability
differences instead of logit gaps), (b)~a lower control penalty that
allows small improvements through, (c)~cumulative flipping (accepting
pairs or triples of groups together to cross the noise threshold), or
(d)~a different search algorithm (evolutionary with population-level
moves, or gradient-free optimization with variance reduction). Naive
greedy hill climbing at group granularity is not viable for this
problem.

% ============================================================
\section{Experiment 12: Corrected Metric Replication}
\label{sec:exp12}

Experiment~12 repeats Experiment~6 with everything held fixed---the same 90
prompts, the same layers $[1,2,3,4,34]$, the same 300 iterations, the same seed,
the same control probes, the same mean fitness---and changes only the
measurement. Search fitness becomes \texttt{seq\_logprob}; the reported metric
becomes greedy generation accuracy, which depends on no distractor token.

The probe set keeps the original prompts and repairs the three defects of
Section~\ref{sec:errata}: full answer strings rather than single token ids,
per-probe distractors reflecting the mistake a model actually makes, and a single
answer convention for trigonometry. Probes may declare alternate renderings, so a
model writing $0.7071$ for $\sqrt{2}/2$ is not marked wrong.

\begin{table}[h]
\centering
\caption{Experiment~12, validation set (30 held-out probes), greedy generation.}
\begin{tabular}{lc}
\toprule
Metric & Count \\
\midrule
Fixed (wrong $\to$ right) & \textbf{5} \\
Broke (right $\to$ wrong) & \textbf{0} \\
Accuracy & 14/30 $\to$ \textbf{19/30} \\
Patch size & 78 flips, 936 bytes \\
\bottomrule
\end{tabular}
\end{table}

All five fixes are changes in what the model emits:
\verb|d/dx [x^7 + x]| from \texttt{0} to \verb|7x^6 + 1|; the second derivative of
$3x^3$ from \verb|18x^2| to \verb|18x|; \texttt{Is 113 prime?} and
\texttt{Is 53 prime?} both from \texttt{No} to \texttt{Yes}; and
\verb|d/dx [e^(-2x)]| from a malformed \verb|-2e^(-2x) + 0 = -2| to
\verb|-2e^(-2x)|. Primality improves most (2/5 to 4/5), consistent with it being
the one original category whose contrast alternated direction and so could not be
won by a token bias. One training-set probe regresses (\texttt{int\_train\_7}) and
is reported rather than screened out.

This is a strictly harder bar than Experiment~6's: the model must emit the
complete correct expression rather than rank one token above another. The
comparison is therefore not $4 \to 2$ but \emph{4 claimed under a broken metric}
versus \emph{5 demonstrated under a sound one}.

% ============================================================
\section{Experiment 13: Corrected Variation Testing}
\label{sec:exp13}

Experiment~13 repeats Experiment~5 as an evaluation only: the same 6-probe
Experiment~4 patch applied to the same 90 novel variation prompts, with the probe
set repaired as in Experiment~12 and scored by full-answer log-probability and
greedy generation.

\begin{table}[h]
\centering
\caption{Experiment~5 as published versus Experiment~13 corrected, 90 variations.}
\begin{tabular}{lcc}
\toprule
Metric & Experiment~5 & Experiment~13 \\
\midrule
Fixed (wrong $\to$ right) & 2 & 2 \\
Broke (right $\to$ wrong) & 5 & \textbf{15} \\
Accuracy & 62/90 $\to$ 59/90 & 49/90 $\to$ \textbf{36/90} \\
\bottomrule
\end{tabular}
\end{table}

Under \texttt{seq\_logprob} the picture is the same: 65/90 to 57/90, with 0 fixed
and 8 broken. Damage appears in five of six categories.

\textbf{Finding.} Experiment~5's conclusion holds and was understated. A patch
trained on 6~probes does not merely fail to generalize; it actively degrades
capability. The original writeup rationalized its 5~breaks as ``all borderline
cases with baseline gaps $< 1.2$''; that explanation does not survive correction,
because the damage is broad rather than marginal.

This is also evidence about the audit itself. The same measurement fix that
\emph{reduced} Experiment~6's claimed benefit \emph{increased} Experiment~5's
measured harm, so the defects were not biased toward flattering results.

Taken together, Experiments~12 and~13 sharpen the paper's central empirical
claim. Training on 6~probes gives 2~fixed and 15~broken; training on 60~diverse
probes gives 5~fixed and 0~broken on held-out prompts. The contrast between
memorization and generalization is starker under sound measurement than under the
original metric.

% ============================================================
\section{Errata and Corrections}
\label{sec:errata}

Experiments~5 and~6 report inflated fix counts. The underlying finding---that
ultra-sparse XOR patches produce targeted behavioral change with no measured
collateral damage---replicates under corrected measurement. Version~1 results
are annotated rather than rewritten, and the published artifact is preserved at
git tag \texttt{v1.0}.

\subsection{Origin}

An independent reproduction in a standalone PyTorch Q1\_0 engine reported two
problems: the hardcoded wrong token was not the model's actual top competitor on
74~of~90 probes, and the integral category was a single token contrast repeated
fifteen times rather than fifteen independent measurements. Both replicate here
(76/90 and confirmed respectively). Auditing them surfaced a third, larger defect
not previously reported. The reproduction also confirmed the baseline of 17/30
validation probes wrong, agreeing across MLX and GGUF/PyTorch runtimes.

\subsection{Defect 1: Dead Probes}

\texttt{encode\_token()} keeps only the last subtoken of an answer string, and
Bonsai's tokenizer splits digits into single characters. Both \texttt{" 20"} and
\texttt{" 0"} therefore reduce to the same id, making
$\text{correct\_id} = \text{wrong\_id}$ and the logit gap \emph{identically zero
regardless of any weight flip}. Such a probe cannot be fixed, cannot be broken,
and scores as wrong because its gap is not positive. Experiment~5 contains 7 such
probes and Experiment~6 contains 8, four of them inside the 30-probe validation
set. The reported baseline of ``17 of 30 wrong'' therefore included 4 probes
incapable of being either right or wrong: the fixable pool was 13, not 17, so the
reported $4/17$ (23.5\%) used an inflated denominator.

\subsection{Defect 2: Measurement Offset}

68 of 90 Experiment~6 answers are multi-token, so the scored id is not the token
the model emits next; on 52 of 90 probes the model's true top-1 next token is a
bare space, and the digit was scored one position early. Under \texttt{token\_gap}
the model appears to answer 13/30 validation probes correctly; scoring the full
answer string gives 22/30 and greedy decoding gives 17/30.

\subsection{Defect 3: Distractor Quality and Category Degeneracy}

The hardcoded distractor was not the model's top competitor on 76/90 probes,
median rank~23. Contrast diversity was uneven across categories: 10 distinct
(correct, wrong) pairs for polynomial and exponential derivatives, 8 for second
derivatives, 3 for trigonometry, 2 for primality, and \textbf{1} for integrals.
The integral category is one measurement reported as fifteen. Primality is a
two-token contrast but alternates direction across probes, so it cannot be won by
a token bias and is not degenerate in the same way.

\subsection{Corrected Result}

Re-scoring the same 93-flip Experiment~6 patch under full-answer log-probability,
changing nothing else, gives 2 held-out fixes rather than 4, with zero breakage.
Three of the four published fixes (\texttt{pd\_val\_1}, \texttt{sd\_val\_3},
\texttt{int\_val\_4}) were already correct at baseline and were counted as
failures only because of Defect~2; \texttt{int\_val\_4} is the integral probe from
the degenerate category. The corrected metric additionally finds one fix the
original missed (\texttt{pd\_val\_0}). Greedy decoding confirms both survivors:
\verb|d/dx [x^7 + x] =| changes from \texttt{0} to \verb|7x^6 + 1|, and
\texttt{Is 113 prime?} changes from \texttt{No} to \texttt{Yes}.

The claim of zero broken probes on held-out data holds under every metric tested.
Experiment~5's sign-flip counts were affected by its 7 dead probes and
understated the harm: rerun as Experiment~13 (Section~\ref{sec:exp13}), the
6-probe patch breaks 15 probes rather than 5. The defects were not biased toward
flattering results---the same fix that reduced Experiment~6's claimed benefit
increased Experiment~5's measured damage.

Correcting the objective did more than deflate the original number. Re-running
the search itself against the sound metric (Experiment~12, Section~\ref{sec:exp12}),
holding layers, iterations, seed and control probes fixed, yields \textbf{5}
held-out fixes with \textbf{zero} breakage from a \emph{smaller} patch---78 flips
and 936~bytes against Experiment~6's 93 flips and 1{,}116~bytes---with every fix
verified in free generation. The original search had spent its budget partly
optimizing a mis-tokenized target scored against a rank-23 distractor; aiming it
correctly recovered more than the defect had cost.

\subsection{Methodological Change}

Logit gaps are retained as a search signal but are no longer reported as a
result. This implements the fix identified by Experiment~11's own
interpretation---a more sensitive fitness function based on log-probability
differences rather than logit gaps. The original metric remains available so
Experiments~1--11 stay reproducible as published.

\section{Limitations}
\label{sec:limitations}

\textbf{Evaluation harness limitations.} Our GSM8K accuracy (22\%) is
well below reported benchmarks, indicating evaluation setup differences.
Logit gap probes don't always predict generation outcomes (visible in
the $7\times8$ example). Standard benchmark evaluation with established
harnesses is a next step.

\textbf{Greedy search finds local optima.} Population-based evolutionary
search with crossover (XOR of XOR patches is a valid patch) could find
better solutions.

\textbf{Row-level granularity is the right operating point for mean-fitness
search.} Experiment~11 confirmed that naive per-group search fails:
individual 128-bit flips produce probe-gap changes an order of magnitude
smaller than row-level flips, and the control-degradation penalty dominates
this tiny signal. Finer granularity requires a more sensitive fitness
function or cumulative flipping (accepting groups of groups together).

\textbf{Attention projections do not help generalization.} Experiment~10
showed that XOR flips on attention Q/K/V/O are mechanically valid but
encode context-specific routing patterns that overfit to training probes
and regress on held-out ones. MLP is the right target for behavioral
patching.

\textbf{GGUF and MLX versions of Bonsai are not the same model.} GGUF
Q1\_0\_g128 is scale-only (1.125~bpw) while MLX is scale+bias (1.25~bpw),
making GGUF strictly less expressive. Identical probe prompts produce
different logit gaps on the two formats, and patches found on one do
not necessarily transfer. Experiments 3--8 use MLX; 9--11 use GGUF.

\textbf{Patch stacking shows interference.} Experiment~7 confirms
stacking is mechanically sound but behaviorally lossy---individual
improvements partially cancel. Joint optimization would likely
outperform naive stacking.

\textbf{Single model.} Bonsai~8B is currently the only production-quality
true 1-bit LLM. The approach does not extend to ternary models.

\textbf{Scale factors are not patched.} Only binary weights are modified,
not FP16 scale/bias values.

\textbf{Small evaluation set.} Our probes cover limited domains.

% ============================================================
\section{Responsible Use}

Bankai modifies model behavior with kilobyte-scale patches invisible at
inference time. The same mechanism could insert subtle malicious changes.
However, XOR patches are transparent by design: every patch is a readable
manifest listing exactly which rows were flipped, and verification is
trivial. If patch libraries become a deployment pattern, provenance and
verification become critical infrastructure.

% ============================================================
\section{Future Work}

\begin{itemize}[nosep]
  \item \textbf{Finer granularity:} Group-level (128~bits) or per-bit
    search for more precise patches.
  \item \textbf{Evolutionary search:} Population-based with crossover
    via XOR and Hamming-distance diversity pressure.
  \item \textbf{Benchmark evaluation:} MMLU subcategories, GSM8K,
    HumanEval.
  \item \textbf{Patch stacking:} Empirical composability testing.
  \item \textbf{Cross-model extraction:} XOR between Bonsai variants
    as naturally occurring patches.
  \item \textbf{Calculus-specific layer selection:} Searching layers~5,
    6, and~10 (identified as high-impact for calculus).
\end{itemize}

% ============================================================
\section{Conclusion}

We have introduced Bankai, the first post-training adaptation method for
true 1-bit LLMs. By exploiting the binary nature of weights, behavioral
modifications reduce to sparse XOR bitmasks that are orders of magnitude
smaller than existing adapters, incur zero inference overhead, and can be
applied and reverted in microseconds.

Through eleven experiments on Bonsai~8B, we characterized the method's
capabilities and limitations: targeted corrections work reliably,
generalization requires diverse training signal, row-level MLP flips
are the right operating point (finer granularity fails at naive greedy
search; attention projections overfit), and the approach ports to
NVIDIA GPUs via a custom llama.cpp tool that runs $\sim$24$\times$ faster
than the Apple Silicon baseline. The method represents a new point in
the adaptation design space---one that is uniquely enabled by true
binary model architectures.

Code, patches, and all experiments are available at
\url{https://github.com/nikshepsvn/bankai}.

% ============================================================
\bibliographystyle{plainnat}
\begin{thebibliography}{12}

\bibitem[Hu et~al.(2021)]{hu2021lora}
E.~J. Hu, Y.~Shen, P.~Wallis, Z.~Allen-Zhu, Y.~Li, S.~Wang, L.~Wang,
  and W.~Chen.
\newblock Lo{RA}: Low-rank adaptation of large language models.
\newblock \emph{arXiv:2106.09685}, 2021.

\bibitem[Rakin et~al.(2019)]{rakin2019bitflip}
A.~S. Rakin, Z.~He, and D.~Fan.
\newblock Bit-flip attack: Crushing neural network with progressive bit
  search.
\newblock In \emph{ICCV}, 2019.

\bibitem[Dong et~al.(2024)]{dong2024stbllm}
P.~Dong, et~al.
\newblock {STBLLM}: Breaking the 1-bit barrier with structured binary
  {LLMs}.
\newblock In \emph{ICLR}, 2025.

\bibitem[Wang et~al.(2023)]{wang2023bitnet}
H.~Wang, S.~Ma, L.~Dong, S.~Huang, H.~Wang, L.~Ma, F.~Yang, R.~Wang,
  Y.~Wu, and F.~Wei.
\newblock Bit{N}et: Scaling 1-bit transformers for large language models.
\newblock \emph{arXiv:2310.11453}, 2023.

\bibitem[Ma et~al.(2025)]{ma2025bitnet}
S.~Ma, H.~Wang, S.~Huang, et~al.
\newblock Bit{N}et b1.58 2{B}4{T} technical report.
\newblock \emph{arXiv:2504.12285}, 2025.

\bibitem[{PrismML}(2026)]{prismml2026}
{PrismML}.
\newblock Bonsai 8{B}.
\newblock \url{https://prismml.com/news/bonsai-8b}, 2026.

\bibitem[Meng et~al.(2022)]{meng2022rome}
K.~Meng, D.~Bau, A.~Andonian, and Y.~Belinkov.
\newblock Locating and editing factual associations in {GPT}.
\newblock In \emph{NeurIPS}, 2022.

\bibitem[Meng et~al.(2023)]{meng2023memit}
K.~Meng, A.~S. Sharma, A.~Andonian, Y.~Belinkov, and D.~Bau.
\newblock Mass-editing memory in a transformer.
\newblock In \emph{ICLR}, 2023.

\bibitem[Bai et~al.(2021)]{bai2021tbfa}
Y.~Bai, et~al.
\newblock Targeted attack against deep neural networks via flipping limited
  weight bits.
\newblock In \emph{ICLR}, 2021.

\bibitem[Xu et~al.(2024)]{xu2024recast}
Y.~Xu, et~al.
\newblock {RECAST}: Reparameterized, compact weight adaptation for
  sequential tasks.
\newblock \emph{arXiv:2411.16870}, 2024.

\bibitem[Bulat and Tzimiropoulos(2020)]{bulat2020xornet}
A.~Bulat and G.~Tzimiropoulos.
\newblock {XOR-Net}: An efficient computation pipeline for binary neural
  network inference on edge devices.
\newblock 2020.

\bibitem[Rastegari et~al.(2016)]{rastegari2016xnor}
M.~Rastegari, V.~Ordonez, J.~Redmon, and A.~Farhadi.
\newblock {XNOR-Net}: Image{N}et classification using binary convolutional
  neural networks.
\newblock In \emph{ECCV}, 2016.

\bibitem[Dettmers et~al.(2023)]{dettmers2023qlora}
T.~Dettmers, A.~Pagnoni, A.~Holtzman, and L.~Zettlemoyer.
\newblock {QLoRA}: Efficient finetuning of quantized language models.
\newblock In \emph{NeurIPS}, 2023.

\end{thebibliography}

\end{document}
